\chapter{Irregular Periods}

\section*{Definition}
Menstrual cycles that vary significantly from the typical 21-35 day pattern, including absent periods (amenorrhea), infrequent periods (oligomenorrhea), or unpredictable cycle length.

\section*{Pathophysiology}
Normal menstruation requires intact hypothalamic-pituitary-ovarian axis with appropriate hormone levels. Disruption at any level causes irregular bleeding: hypothalamic (stress, exercise, weight loss), pituitary (prolactinoma), ovarian (PCOS, perimenopause, premature ovarian failure), uterine (fibroids, polyps), or endocrine (thyroid, adrenal) causes.

\section*{Etiology}
\begin{itemize}
  \item Polycystic ovary syndrome (PCOS)
  \item Perimenopause (transition to menopause)
  \item Thyroid disorders (hyperthyroidism or hypothyroidism)
  \item Hyperprolactinemia
  \item Stress or significant weight changes
  \item Excessive exercise (athletic amenorrhea)
  \item Eating disorders (anorexia nervosa)
  \item Uterine fibroids or polyps
  \item Primary ovarian insufficiency (premature menopause)
  \item Hormonal contraception effects
\end{itemize}

\section*{Indications for Evaluation}
Seek care if:
\begin{itemize}
  \item Severe or worsening symptoms
  \item Missed periods for 3+ cycles, very heavy bleeding, bleeding between periods, or difficulty conceiving
  \item Accompanied by other concerning signs
\end{itemize}

\section*{Related Symptoms}
\begin{itemize}
  \item Heavy bleeding
  \item Pelvic pain
  \item Weight changes
  \item Hair growth changes
  \item Acne
\end{itemize}
\textbf{Note:} Always consider serious underlying conditions when evaluating persistent irregular periods.

\section*{Self-Care / Home Management}
\begin{itemize}
  \item Rest and avoid activities that may aggravate the condition.
  \item Maintain adequate hydration and a balanced diet.
  \item Over-the-counter remedies may provide symptomatic relief (consult a pharmacist or doctor).
  \item Monitor symptoms and seek professional advice if they persist or worsen.
  \item Avoid self-medicating with prescription medications without medical guidance.
\end{itemize}

\section*{Diagnostic Approach}
\begin{itemize}
  \item A thorough history and physical examination are the first steps in evaluation.
  \item Laboratory tests (blood work, urinalysis) may be ordered based on clinical suspicion.
  \item Imaging studies (X-ray, ultrasound, CT, MRI) may be indicated when structural pathology is suspected.
  \item Specialist referral is arranged when the diagnosis remains uncertain or requires specific expertise.
\end{itemize}

\section*{Treatment / Management Overview}
\begin{itemize}
  \item Treatment targets the underlying cause identified through diagnostic evaluation.
  \item Supportive care and symptom management are provided while awaiting definitive therapy.
  \item Medications, lifestyle modifications, and physical therapies may be prescribed as appropriate.
  \item Follow-up ensures treatment effectiveness and allows timely adjustments to the care plan.
\end{itemize}

\clearpage